\lysh{20647}{4}
\begin{alist}
\item Αν υποθέσουμε ότι το πολυώνυμο έχει ακέραιους συντελεστές, τότε κάθε ακέραια ρίζα του θα είναι διαιρέτης του σταθερού όρου του. Ο σταθερός όρος του όμως είναι ίσος με 3 και ο αριθμός 2 που είναι ρίζα του πολυωνύμου δεν είναι διαιρέτης του 3. Άρα το πολυώνυμο δεν έχει όλους τους συντελεστές του ακέραιους, οπότε υπάρχει ένας τουλάχιστον συντελεστής του που δεν είναι ακέραιος.
\item Με την επιπλέον πληροφορία ότι $P(1) = 0$ έχουμε:
\begin{align*}
\begin{cases} P(2) = 0 \\ P(1) = 0 \end{cases} &\Leftrightarrow \begin{cases} 8\alpha + 4\beta - 2\beta + 3 = 0 \\ \alpha + \beta - \beta + 3 = 0 \end{cases} \\
&\Leftrightarrow \begin{cases} \alpha = -3 \\ 2\beta = 21 \end{cases} \Leftrightarrow \begin{cases} \alpha = -3 \\ \beta = \dfrac{21}{2} \end{cases}
\end{align*}
\item Με $\alpha = -3$ και $\beta = \dfrac{21}{2}$ έχουμε $P(x) = -3x^3 + \dfrac{21}{2}x^2 - \dfrac{21}{2}x + 3$. \\
Αρχικά θα βρούμε τις ρίζες του πολυωνύμου. \\
Είναι: 
\[P(x) = 0 \Leftrightarrow -3x^3 + \frac{21}{2}x^2 - \frac{21}{2}x + 3 = 0 \Leftrightarrow -6x^3 + 21x^2 - 21x + 6 = 0\]
Είναι όμως γνωστό ότι $P(1) = 0$, οπότε με τη βοήθεια του διπλανού σχήματος Horner έχουμε:
\begin{center}
\Horner{6,21,-21,6}{1}
\end{center}
\begin{align*}
P(x) = 0 &\Leftrightarrow (x-1)(-6x^2 + 15x - 6) = 0 \\
&\Leftrightarrow x - 1 = 0 \text{ ή } -6x^2 + 15x - 6 = 0
\end{align*}
Η εξίσωση $-6x^2 + 15x - 6 = 0$ που είναι ισοδύναμη με την $2x^2 - 5x + 2 = 0$ έχει ρίζες τους αριθμούς $2$ και $\dfrac{1}{2}$. Επομένως το πολυώνυμο έχει ρίζες τους αριθμούς $\dfrac{1}{2}, 1$ και $2$. \\
Έχουμε λοιπόν: 
\[P(x) \le 0 \Leftrightarrow -6(x-1)(x-2)\left(x-\frac{1}{2}\right) \le 0\]
και όταν $x>2$ όλοι οι παράγοντες του πολυωνύμου είναι θετικοί, οπότε αυτό είναι ομόσημο του $-6$, δηλαδή αρνητικό. Σε καθένα από τα προηγούμενα διαστήματα αλλάζει πρόσημο ακριβώς ένας όρος του γινομένου, οπότε προκύπτει ο παρακάτω πίνακας προσήμων του πολυωνύμου.
\begin{center}
\begin{tabular}{|c|c|c|c|c|c|c|c|}

$x$ & $-\infty \hspace{2cm} \dfrac{1}{2}$ & $1$ & $2$ & $+\infty$ \\ 
$P(x)$ & $+$ \hspace{0.5cm} $0$ \hspace{0.5cm} $-$ & $0$ \hspace{0.5cm} $+$ & $0$ \hspace{0.5cm} $-$ \\ 
\end{tabular}
\end{center}
\end{alist}

